\PassOptionsToPackage{hyphens}{url}
\documentclass[sigconf,screen,review,anonymous]{acmart}

\usepackage{xcolor}
\usepackage{xurl}
\usepackage{microtype}
\usepackage{tabularx}
\usepackage{array}
\usepackage{makecell}
\usepackage{ragged2e}

\makeatletter
\AtBeginDocument{%
  \providecommand\BibTeX{{%
    Bib\TeX}}%
  \@ifundefined{linenumberfont}{}{%
    \renewcommand\linenumberfont{\normalfont\tiny\sffamily\color[gray]{0.55}}}}
\makeatother

\newcolumntype{L}[1]{>{\RaggedRight\arraybackslash}p{#1}}
\newcolumntype{C}[1]{>{\Centering\arraybackslash}p{#1}}
\newcolumntype{Y}{>{\RaggedRight\arraybackslash}X}
\renewcommand\theadfont{\small\bfseries}
\setcellgapes{2pt}
\makegapedcells

\setcopyright{none}
\settopmatter{printfolios=true}
\renewcommand\footnotetextcopyrightpermission[1]{}

\acmConference[MM '26]{Proceedings of the 34th ACM International Conference on Multimedia}{November 10--14, 2026}{Rio de Janeiro, Brazil}
\acmBooktitle{Proceedings of the 34th ACM International Conference on Multimedia (MM '26), November 10--14, 2026, Rio de Janeiro, Brazil}
\acmDOI{}
\acmISBN{}

\title{RIIC-ReID: Retrieval-Informed Illumination Correction for Cross-Light Re-Identification}

\author{Anonymous Author(s)}
\affiliation{%
  \institution{Anonymous Institution}
  \city{Anonymous City}
  \country{Anonymous Country}}

\begin{abstract}
Most illumination correction methods optimize images for human observers.
They improve contrast, brightness, and visual naturalness in image space.
Retrieval systems, however, are judged by a different observer: the
downstream model.
This creates a structural mismatch between human-perceptual enhancement and
model-centric adaptation.
We study this mismatch as a general problem in re-identification and
ground it in wildlife re-identification, where cross-camera illumination
shift is especially severe and identity often depends on local texture cues.
We therefore formulate illumination handling not as standalone enhancement,
but as constrained, retrieval-driven adaptation.
Weak photometric priors define a feasible correction set, while retrieval
geometry selects the preferred correction inside that set.
We instantiate this view as RIIC-ReID, a retrieval-informed illumination
correction framework that couples a compact illumination operator to the re-identification
backbone through three principles: feasible correction, trust-controlled
adaptation, and geometry-guided supervision.
The framework keeps raw, safe base-corrected, and refined branches, and uses
learnable trust gates to select among them per local region.
On ATRW under a matched protocol, RIIC-ReID raises mmAP from a baseline of
$56.69$ to $58.11$ and remains stronger than an illumination-only front-end.
Standardized fixed query-gallery experiments on GZGC Zebra and StripeSpotter
show the same direction, while generic perceptual enhancers such as
Zero-DCE++ and RetinexNet remain weaker or less consistent.
These results support a simple conclusion: for retrieval, illumination
correction should be optimized for the model rather than for human perception
alone.
\end{abstract}

\begin{CCSXML}
<ccs2012>
 <concept>
  <concept_id>10010147.10010257.10010293</concept_id>
  <concept_desc>Computing methodologies~Image representations</concept_desc>
  <concept_significance>500</concept_significance>
 </concept>
 <concept>
  <concept_id>10010147.10010257.10010282.10010284</concept_id>
  <concept_desc>Computing methodologies~Image processing</concept_desc>
  <concept_significance>300</concept_significance>
 </concept>
 <concept>
  <concept_id>10002951.10003317.10003347.10003350</concept_id>
  <concept_desc>Information systems~Retrieval models and ranking</concept_desc>
  <concept_significance>500</concept_significance>
 </concept>
</ccs2012>
\end{CCSXML}

\ccsdesc[500]{Computing methodologies~Image representations}
\ccsdesc[300]{Computing methodologies~Image processing}
\ccsdesc[500]{Information systems~Retrieval models and ranking}

\keywords{illumination correction, re-identification, retrieval, embedding geometry, task-aligned front-end}

\begin{document}

\maketitle

\section{Introduction}

Most illumination correction methods are built for human observers.
They improve contrast, brightness, and visual naturalness in image
space~\cite{land1971lightness,land1977retinex,wei2018retinexnet,guo2020zerodce,jiang2021enlightengan,wu2022uretinex,ma2022sci,xu2022snraware,cai2023retinexformer,fu2023pairlie,xu2023structureguidance,yi2023diffretinex,yang2023nerco,liu2023msrq,wang2024quadprior,shi2024zeroig,saini2024specularity}.
That objective is sensible when the final consumer is a person, but in
retrieval systems the final observer is not the human eye: it is the
downstream model.
This raises a more fundamental question than low-light enhancement itself:
\emph{should illumination correction still be defined by human perception, or
should it adapt to the feature space preferred by the model?}

This question is general across re-identification, but it becomes especially
visible in wildlife settings, where capture conditions are less controlled and
identity often depends on fine-grained local texture and species-specific
pattern cues~\cite{schneider2019past,schneider2022similarity,cermak2024wildlifedatasets,hou2025openanimals}.
In such cases, even strong modern encoders must absorb large camera- and
illumination-induced appearance shifts~\cite{luo2019bag,zhou2019osnet,he2021transreid,zhou2021osnetain,fu2022pnl,dou2023isr,chen2023solider,zhu2024seas,ni2023partaware,he2024instructreid}.
Two images of the same individual captured under different cameras or lighting
conditions can differ more in brightness, local shadow, highlight pattern, and
color cast than two images of different individuals captured under similar
lighting.
Once this happens, illumination handling is no longer a cosmetic preprocessing
detail; it directly reshapes the feature geometry on which retrieval succeeds
or fails.
An image that looks more natural to a person need not be the image that best
preserves identity-discriminative structure for a model
(Figure~\ref{fig:retrieval-comparison}).

This paper studies that mismatch explicitly.
Our motivation is broad, but the empirical focus of this paper is
cross-light re-identification.
We do not propose a new physical theory of illumination; rather,
we argue that illumination correction should be \emph{observer-dependent}:
human-centered enhancement and model-centered adaptation are related but
different optimization problems.
Retinex-style structure~\cite{land1971lightness,land1977retinex} remains useful,
but only as a weak admissibility prior~\cite{wei2018retinexnet,wu2022uretinex,xu2022snraware,cai2023retinexformer,yi2023diffretinex,wang2024quadprior}.
Once this distinction is made explicit, illumination handling should no longer
be treated as an isolated enhancement block placed before recognition.
It should instead be treated as a constrained and learnable front-end of the
retrieval system itself~\cite{lengyel2021physicsprior,liu2022imageadaptive,hashmi2023featenhancer,wu2023semanticaware,luo2023simminmax,du2024daynight,wang2024motdark}.
In this view, photometric plausibility defines what
corrections are admissible, and retrieval behavior determines which admissible
correction is actually preferred.
We validate this view primarily in wildlife re-identification, where
illumination variation is a first-order nuisance and local texture makes
correction errors especially costly; we use that setting as a stress-test
domain rather than as a claim of universal retrieval coverage.

The challenge is that a purely task-driven corrector is ill-posed.
This tension is already visible in ReID settings that either push
illumination robustness into the descriptor~\cite{huang2019illumination,xu2020blackreid,yang2020contour,xiang2022rethinking,li2024dnvehicle}
or introduce dedicated nighttime enhancement pipelines~\cite{lu2024illuminationdistillation,lu2025nighttime}.
If trained only through retrieval loss, the corrector may discover non-physical or
dataset-specific transformations that reduce training error while harming
generalization.
If constrained too strongly by perceptual priors, it may never move samples
toward the feature structure actually preferred by the encoder.
The problem is therefore not simply to learn a stronger correction function,
but to balance \emph{task alignment}, \emph{feasibility}, and \emph{trust}.
Viewed this way, the value of the method is not a new photometric law; it is
an initial but concrete step toward model-adaptive illumination correction, in
which the front-end is trained for the downstream model rather than for human
perception alone.

We address this problem with RIIC-ReID, a retrieval-informed illumination correction framework
learned jointly with the downstream encoder.
The framework is organized around three principles derived from the task formulation:
\emph{feasible correction}, which bounds the correction inside a safe
photometric set;
\emph{trust-controlled adaptation}, which permits correction only where
image-space and feature-space evidence justify it;
and \emph{geometry-guided supervision}, which lets retrieval structure rather
than perceptual preference determine the final correction direction.

\noindent\textbf{Contributions.}
\textbf{(1)} We introduce a general model-centric view of illumination
correction for retrieval, where correction should follow the downstream
model's discriminative geometry and photometric priors act only as
feasibility constraints.
\textbf{(2)} We instantiate this view in RIIC-ReID, a constrained retrieval
front-end that combines bounded inverse illumination, model-aware residual
adaptation, identity protection, stripe-wise rollback, multi-branch fusion,
and teacher-guided geometry. We study it primarily in wildlife
re-identification, where illumination variation and texture-dependent identity
cues are tightly coupled, making correction errors especially costly and the
setting a natural stress test for feasibility and trust mechanisms.
\textbf{(3)} We provide matched empirical evidence in that setting: on ATRW,
naive illumination drops mmAP from the baseline's $56.69$ to $55.62$, while full RIIC-ReID
reaches $58.11$; the same ordering remains visible across backbones,
standardized fixed query-gallery tests on GZGC Zebra and StripeSpotter, and
perceptual baselines such as Zero-DCE++ and RetinexNet.

\begin{figure*}[t]
  \centering
  \includegraphics[width=0.985\textwidth]{figures_generated/fig1_retrieval.pdf}
  \caption{Retrieval comparison on a representative ATRW query (ID~231,
  luminance 0.57).
  \textbf{(a)}~Top-5 gallery results under four illumination-handling
  strategies.  Green borders denote correct matches; red borders denote wrong
  identities.  RetinexNet and Zero-DCE++ visually improve the query but
  introduce confuser retrievals (ID~232), while RIIC-ReID returns all correct
  matches.
  \textbf{(b)}~Correction comparison: the four query variants differ subtly
  in image space, yet produce markedly different retrieval outcomes.
  \textbf{(c)}~Average cosine similarity between the query embedding and
  same-identity vs.\ confuser gallery embeddings.  Perceptual methods invert
  the margin (query becomes closer to the confuser), whereas RIIC-ReID
  achieves the largest positive margin ($\Delta{=}+0.019$).}
  \Description{A three-part qualitative figure comparing one ATRW tiger query
  under raw input, RetinexNet, Zero-DCE++, and RIIC-ReID. The left panel shows
  top-5 retrieval lists with correct matches highlighted in green and confusers
  in red. The middle panel shows the four query appearances after different
  illumination treatments. The right panel plots cosine similarity to the same
  identity and to a confuser, showing that RIIC-ReID preserves a positive
  retrieval margin while perceptual enhancement can invert it.}
  \label{fig:retrieval-comparison}
\end{figure*}

\section{Related Work}

\subsection{Photometric and Retinex-Style Correction}

Classical Retinex theory~\cite{land1971lightness,land1977retinex} models an image as the
product of reflectance and illumination, then recovers a reflectance-like
representation by dividing out the estimated illumination.
Deep Retinex methods, curve-estimation models, and recent unfolding, SNR-aware,
structure-guided, transformer, diffusion, and implicit-representation variants
learn this decomposition or related image-formation constraints end-to-end~\cite{wei2018retinexnet,guo2020zerodce,jiang2021enlightengan,wu2022uretinex,ma2022sci,xu2022snraware,fu2023pairlie,xu2023structureguidance,cai2023retinexformer,yi2023diffretinex,yang2023nerco,liu2023msrq,wang2024quadprior,shi2024zeroig,saini2024specularity}.
These methods can substantially improve visual quality, but their target remains
photometric plausibility rather than retrieval utility.

\subsection{Task-Aware Front-Ends for Recognition}

Another line of work evaluates front-end processing through downstream utility
rather than perceptual fidelity alone.
In adverse-weather, day-night, and low-light recognition, task-conditioned
front-ends adapt image processing to detection or tracking heads and improve
downstream performance~\cite{lengyel2021physicsprior,liu2022imageadaptive,hashmi2023featenhancer,wu2023semanticaware,luo2023simminmax,du2024daynight,wang2024motdark}.
Related work also studies day-night transfer for retrieval itself~\cite{li2024dnvehicle}.
Our work belongs to this task-aware family, but differs in making the mismatch
between perceptual and retrieval optima explicit and using it to design both
the front-end and its supervision.

\subsection{Re-Identification Under Illumination Shift}

Re-identification methods often absorb illumination variation indirectly
through augmentation, stronger evaluation protocols, large-scale pre-training,
and modern CNN/Transformer backbones~\cite{zheng2015market,ristani2016performance,luo2019bag,zhou2019osnet,he2021transreid,zhou2021osnetain,fu2022pnl,dou2023isr,chen2023solider,zhu2024seas,ni2023partaware,he2024instructreid}.
Several works address illumination more explicitly.
Illumination-invariant ReID and Black Re-ID~\cite{huang2019illumination,xu2020blackreid}
design representations intended to be more stable under difficult lighting,
while later work studies domain generalization, relevance-aware experts,
person-scale meta-generalization, exemplar-memory adaptation, clothing change,
visible-infrared transfer, and explicit nighttime matching~\cite{zhao2021m3l,dai2021ramoe,bai2021person30k,zhong2019invariance,yang2020contour,ye2020ddag,wei2020coal,yang2022adca,xiang2022rethinking,jiang2025privacyvir,lu2024illuminationdistillation,lu2025nighttime}.
Wildlife and animal re-identification has likewise expanded from tiger-specific
pipelines to broader surveys, ecological analyses, toolkits, and multi-species
resources~\cite{freytag2016chimpanzee,liu2019pgcfl,liu2019ppgnet,li2020atrw,schneider2019past,schneider2022similarity,cermak2024wildlifedatasets,adam2024seaturtleid,nepovinnykh2024species,wahltinez2024opensource,moosa2024selfsupervised,hou2025openanimals,adam2025wildlifereid10k}.
Our view is complementary: instead of pursuing illumination robustness only in
the descriptor or the dataset, we learn a bounded front-end whose correction is
selected by downstream retrieval behavior.
In this paper, we study that idea through a retrieval-centered formulation and
use wildlife benchmarks as a stress-test domain where illumination errors are
especially exposed.

\section{Why Perceptual Illumination Correction Is Not Retrieval-Optimal}

\subsection{Task-Oriented Formulation}

Let $x \in [0,1]^{3 \times H \times W}$ denote an RGB image, $y$ its identity
label, $g_{\phi}$ a learnable illumination front-end, and $f_{\theta}$ the
downstream retrieval encoder.
We denote by
\begin{equation}
  \tilde{x} = g_{\phi}(x), \qquad z = f_{\theta}(\tilde{x})
\end{equation}
the corrected image and its embedding, respectively.

For retrieval-oriented illumination handling, we solve
\begin{equation}
  \min_{\phi,\theta}
  \mathbb{E}_{(x,y)} \left[ \mathcal{L}_{\text{task}}(f_{\theta}(g_{\phi}(x)), y) \right]
  \quad
  \text{s.t.}
  \quad
  g_{\phi}(x) \in \mathcal{F}(x),
  \label{eq:main-problem}
\end{equation}
where $\mathcal{F}(x)$ is a feasible set that enforces bounded gain, bounded
chromatic drift, and identity-preserving deviations from the original image.

Equation~\eqref{eq:main-problem} differs from conventional enhancement
objectives~\cite{wei2018retinexnet,guo2020zerodce,jiang2021enlightengan,cai2023retinexformer,yi2023diffretinex},
in that visual plausibility is enforced only through the feasible set
$\mathcal{F}(x)$; the optimization objective itself is entirely
retrieval-driven.

\subsection{Why the Two Optima Generally Differ}

To make the mismatch explicit, define the retrieval objective
$L_x(\tilde{x}) = \mathcal{L}_{\text{task}}(f_{\theta}(\tilde{x}), y)$
and a photometric objective $P_x(\tilde{x})$ specified without reference to
downstream retrieval geometry.
Within the same feasible set $\mathcal{F}(x)$, consider three solutions:
\[
\begin{aligned}
  \tilde{x}_P^{\star}
  &=
  \arg\min_{\tilde{x} \in \mathcal{F}(x)} P_x(\tilde{x}), \\
  \tilde{x}_T^{\star}
  &=
  \arg\min_{\tilde{x} \in \mathcal{F}(x)} L_x(\tilde{x}), \\
  \tilde{x}_{\lambda}^{\star}
  &=
  \arg\min_{\tilde{x} \in \mathcal{F}(x)}
  \bigl(L_x(\tilde{x}) + \lambda P_x(\tilde{x})\bigr).
\end{aligned}
\]
The first is the photometric optimum, the second the retrieval optimum, and
the third the weak-prior solution used in our formulation.

The non-identifiability of Retinex-style decomposition already implies that
some prior is unavoidable~\cite{land1971lightness,wei2018retinexnet,cai2023retinexformer}:
for any positive field $c$, one may write
$x = (R \odot c) \odot (L / c)$.
The relevant question is therefore not whether a prior exists, but whether it
should dominate the objective or merely regularize the feasible set.

\paragraph{Prop. 1 (Local mismatch).}
Assume $\mathcal{F}(x)$ is an embedded $C^1$ feasible submanifold in a neighborhood of a
photometric optimum $\tilde{x}_P^{\star}$, and $L_x,P_x$ are differentiable on
that neighborhood.
Let
\[
  \mathcal{T}_P = T_{\mathcal{F}(x)}(\tilde{x}_P^{\star}),
  \qquad
  g_P = \Pi_{\mathcal{T}_P}\nabla L_x(\tilde{x}_P^{\star}).
\]
If $g_P \neq 0$, then $\tilde{x}_P^{\star}$ is not a local minimizer of
$L_x$ over $\mathcal{F}(x)$ and therefore cannot coincide with
$\tilde{x}_T^{\star}$.

\paragraph{Proof sketch (full proof in supplementary material).}
Since $g_P \in \mathcal{T}_P$, there exists a feasible $C^1$ curve
$\gamma(t) \subset \mathcal{F}(x)$ with $\gamma(0)=\tilde{x}_P^{\star}$ and
$\gamma'(0)=-g_P$, so
\[
  \frac{d}{dt}L_x(\gamma(t))\Big|_{t=0}
  =
  \langle \nabla L_x(\tilde{x}_P^{\star}), -g_P \rangle
  =
  -\|g_P\|_2^2 < 0.
\]
Hence $L_x$ strictly decreases along a feasible direction through
$\tilde{x}_P^{\star}$, so the photometric optimum cannot be retrieval-optimal on
$\mathcal{F}(x)$ unless the projected retrieval gradient vanishes.

Proposition~1 formalizes the key point of the paper.
A perceptually reasonable correction is not retrieval-optimal unless the
photometric optimum is also stationary for retrieval on the same feasible
manifold.
Such alignment is possible, but it is not guaranteed by Retinex-style
admissibility alone. Conversely, the vanishing of the projected retrieval
gradient is only a necessary local condition, not a sufficient one for the two
optima to coincide.

\subsection{A Local Justification for Weak Priors}

We now justify why the photometric term should be weak.
Introduce local coordinates $\delta$ on $\mathcal{F}(x)$ around a retrieval
optimum $\tilde{x}_T^{\star}$, so $\delta=0$ corresponds to
$\tilde{x}_T^{\star}$.
Assume $L_x(\delta)$ and $P_x(\delta)$ are twice continuously differentiable
and that the restricted Hessian
$H = \nabla^2 L_x(0)$ is positive definite on the feasible tangent space.
Consider the weak-prior objective
\[
  J_{\lambda}(\delta) = L_x(\delta) + \lambda P_x(\delta),
  \qquad
  \lambda \ge 0.
\]

\paragraph{Prop. 2 (Weak-prior perturbation).}
Under the assumptions above, for sufficiently small $\lambda$ there exists a
local minimizer $\delta_{\lambda}^{\star}$ of $J_{\lambda}$ satisfying
\[
  \delta_{\lambda}^{\star}
  =
  -\lambda H^{-1}\nabla P_x(0) + o(\lambda).
\]
Consequently,
\[
  \|\delta_{\lambda}^{\star}\|_2 = O(\lambda),
  \qquad
  L_x(\delta_{\lambda}^{\star}) - L_x(0) = O(\lambda^2).
\]

\paragraph{Proof sketch (full proof in supplementary material).}
The stationarity condition is
\[
  \nabla L_x(\delta_{\lambda}^{\star})
  +
  \lambda \nabla P_x(\delta_{\lambda}^{\star}) = 0.
\]
Because $\delta=0$ is a retrieval optimum, $\nabla L_x(0)=0$.
Since $H$ is invertible, the implicit function theorem gives a $C^1$ local
solution branch and therefore
\[
  \delta_{\lambda}^{\star}
  =
  -\lambda H^{-1}\nabla P_x(0) + o(\lambda).
\]
In particular, $\|\delta_{\lambda}^{\star}\|_2 = O(\lambda)$.
Substituting this into the second-order expansion of $L_x$ around $0$,
\[
  L_x(\delta)
  =
  L_x(0) + \frac{1}{2}\delta^{\top}H\delta + o(\|\delta\|_2^2),
\]
yields the claimed $O(\lambda^2)$ excess retrieval risk.

Proposition~2 shows that a weak photometric prior perturbs the retrieval
solution only to first order, while the induced retrieval risk increases only
to second order.
(The effective range of this regime depends on the condition number of~$H$:
an ill-conditioned Hessian amplifies $\|\delta_{\lambda}^{\star}\|_2$ for a
given~$\lambda$, which motivates the conservative annealing schedule used in
practice.)
This is exactly the regime we want: photometric structure regularizes the local
search space without dictating the correction target.

\subsection{Design Implications}

A final property completes the analysis: once model-aware adaptation moves
beyond bounded photometric correction, the resulting change must remain
trust-bounded.

\paragraph{Prop. 3 (Trust-bounded adaptation).}
Suppose the corrected image is produced by
\[
  \tilde{x}
  =
  \alpha \odot \tilde{x}_{\mathrm{prot}} + (1-\alpha)\odot x,
  \qquad
  0 \le \alpha \le 1,
\]
and a stripe feature is fused by
\[
  F_{\mathrm{fuse}}^{(s)} = \sum_{b=1}^{B}\alpha_b^{(s)}F_b^{(s)},
  \qquad
  \alpha_b^{(s)} \ge 0,\quad \sum_{b=1}^{B}\alpha_b^{(s)}=1.
\]
Then, for any monotone norm,
\[
  \|\tilde{x}-x\|
  \le
  \|\tilde{x}_{\mathrm{prot}}-x\|,
\]
and, for each stripe $s$ and any norm,
\[
  \|F_{\mathrm{fuse}}^{(s)} - F_{\mathrm{raw}}^{(s)}\|
  \le
  \sum_{b \neq \mathrm{raw}}
  \alpha_b^{(s)}
  \|F_b^{(s)} - F_{\mathrm{raw}}^{(s)}\|.
\]

\paragraph{Proof sketch (full proof in supplementary material).}
The image-space bound follows because
$\tilde{x}-x=\alpha\odot(\tilde{x}_{\mathrm{prot}}-x)$ with
$0 \le \alpha \le 1$ elementwise.
The feature-space bound follows because softmax fusion makes each fused stripe a
convex combination of branch features, so
\[
  F_{\mathrm{fuse}}^{(s)} - F_{\mathrm{raw}}^{(s)}
  =
  \sum_{b \neq \mathrm{raw}}
  \alpha_b^{(s)}\bigl(F_b^{(s)} - F_{\mathrm{raw}}^{(s)}\bigr),
\]
 after which the triangle inequality gives the claim.

Proposition~3 explains why rollback and branch fusion are not heuristic extras:
they implement explicit image-space and feature-space trust bounds.

The three propositions now yield a direct algorithmic reading.
Proposition~1 says that photometric admissibility alone cannot determine the
retrieval-preferred correction.
Proposition~2 says that the photometric term should regularize the search only
weakly.
Proposition~3 says that the correction must remain trust-bounded in both image
space and feature space.
These observations lead to three requirements for a retrieval-oriented
illumination front-end.
First, correction must remain \emph{feasible}: the model needs bounded gain,
bounded chromatic drift, and weak physical regularity so the learned operator
does not collapse into arbitrary image manipulation.
Second, correction must be \emph{trust-controlled}: the exact projection onto
$\mathcal{F}(x)$ is intractable once semantic constraints are introduced, so
the system needs learnable image-space and feature-space damping mechanisms.
Third, correction must be \emph{geometry-guided}: once the feasible set is
defined, the preferred correction should be selected by retrieval geometry
rather than by perceptual preference.
RIIC-ReID is a direct implementation of these three requirements.

\section{RIIC-ReID}

\begin{figure*}[t]
  \centering
  \includegraphics[width=0.78\textwidth]{figures_generated/figure2_2.jpg}
  \caption{RIIC-ReID architecture.  A bounded base corrector defines feasible
  illumination correction (\S\ref{sec:feasible}).  A model-aware residual
  branch adapts correction using mid-level encoder features, while identity
  protection and stripe-wise rollback suppress risky edits
  (\S\ref{sec:trust}).  Stripe-aware competitive fusion selects among raw,
  base-corrected, and refined branches.  Teacher-guided geometry losses
  supervise the corrected features to remain within same-identity manifold
  tubes (\S\ref{sec:geometry}).}
  \Description{A schematic of the RIIC-ReID pipeline. The top branch estimates
  luminance and feasible base correction. A middle branch performs
  trust-controlled residual adaptation using mid-level features, identity
  protection, and stripe-wise rollback. The right side fuses raw,
  base-corrected, and refined branches with stripe-aware attention before the
  final descriptor head. The lower panel summarizes the loss terms used for
  retrieval, photometric feasibility, and teacher-guided geometry.}
  \label{fig:architecture}
\end{figure*}

\subsection{Overview}
\label{sec:overview}

Figure~\ref{fig:architecture} gives an overview.
The RIIC-ReID framework couples a compact illumination operator to the retrieval encoder without
turning the system into a free-form enhancer.
The illumination branch first constructs a feasible base correction, then
applies model-aware adaptation and lightweight refinement under explicit trust
control.
The encoder simultaneously keeps a raw branch, so the final descriptor can be
selected from raw, bounded base-corrected, and refined evidence rather than
being forced to trust one corrected image.
This organization mirrors the theory above and stays close to the broader
idea of task-aware front-ends for recognition~\cite{liu2022imageadaptive,hashmi2023featenhancer,wu2023semanticaware}:
photometric priors delimit the
admissible set, trust mechanisms keep adaptation conservative, and retrieval
geometry determines which admissible correction is actually worth using.

\subsection{Feasible Illumination Correction}
\label{sec:feasible}

We adopt the standard reflectance--illumination factorization as an approximate
inductive bias~\cite{land1971lightness,land1977retinex,wei2018retinexnet,cai2023retinexformer},
\begin{equation}
  x \approx R \odot L,
\end{equation}
where $R$ is a reflectance-like term and $L$ is an illumination field.
The factorization is not treated as an exact physical claim; it simply provides
structure for a learnable front-end.

The framework first computes a luminance map
\begin{equation}
  Y = 0.299 x_R + 0.587 x_G + 0.114 x_B,
\end{equation}
then predicts a coarse illumination field $L_0 = E_{\phi}(Y)$.
To allow spatially varying task sensitivity, it also predicts a multiplicative
sensitivity map $S_{\phi}(x)$ and forms
\begin{equation}
  L_c = 1 + (L_0 - 1) \odot S_{\phi}(x).
\end{equation}
This yields a stable low-frequency estimate without forcing the same correction
strength everywhere.

To handle color shift, the front-end can further augment the scalar illumination with a
learned white-balance gain or a learned multi-channel illumination field.
We denote the resulting effective correction field by $\hat{L}$.
Direct inverse illumination correction is unstable in dark or saturated regions,
so the framework applies a bounded scaling rule.
We compute the desired inverse scale
$a = 1/(\hat{L} + \varepsilon)$
and the per-pixel saturation-safe upper bound
$b = 0.99/(x + \varepsilon)$.
The actual scale is obtained through a differentiable soft minimum
$s = \operatorname{softmin}(a,b)$,
then clamped to a bounded gain range together with a bound on chromatic drift
$\|\Delta_{\text{chroma}}\|_{\infty} \le \tau_c$.
The resulting base-corrected image is
$\tilde{x}_{\text{base}} = x \odot s$.
This stage implements the feasible-set part of the theory: it gives the model a
learnable but explicitly bounded correction space.
These constraints are not the final target; they only ensure that subsequent
task-driven adaptation starts from a safe operating region.

\subsection{Trust-Controlled Adaptation}
\label{sec:trust}

Feasible correction alone is not enough, because even an admissible correction
need not be trustworthy everywhere.
The framework therefore introduces two layers of trust control.

First, a safe coarse guidance image is fed into the early layers of the
backbone to extract a mid-level feature map $F_{\text{mid}}$.
This feature map conditions a model-aware residual branch:
\begin{equation}
  \tilde{x}_{\text{att}}
  =
  M_{\phi}(x, \tilde{x}_{\text{base}}, \hat{L}, F_{\text{mid}}).
\end{equation}
The correction therefore responds to the encoder's internal representation
rather than to image statistics alone.
This follows prior task-aware enhancement pipelines that use downstream
semantics to guide low-level processing~\cite{liu2022imageadaptive,hashmi2023featenhancer,wu2023semanticaware},
but here the signal is explicitly retrieval-driven.

Second, the model restricts where and how strongly this residual may act.
An identity protection map
$P_{\text{id}} \in [0,1]^{1 \times H \times W}$ is constructed from early
backbone activations and used to suppress changes in identity-critical regions:
\begin{equation}
  \tilde{x}_{\text{prot}}
  =
  x + (1 - P_{\text{id}}) \odot (\tilde{x}_{\text{att}} - x).
\end{equation}
It then predicts a stripe-wise rollback gate $\alpha$, conditioned on
illumination, color risk, correction magnitude, and $F_{\text{mid}}$, so each
horizontal body region can independently return toward the raw observation:
\begin{equation}
  \tilde{x}
  =
  \alpha \odot \tilde{x}_{\text{prot}} + (1-\alpha) \odot x.
\end{equation}
This image-space rollback acts as a learned trust-region mechanism: risky
regions receive smaller effective correction steps.
A lightweight refiner then sharpens local detail around the rollback solution,
and the identity protection mask is applied again so refinement cannot
overwrite protected regions.

Trust control continues in feature space.
The framework preserves a raw branch in the encoder and fuses raw, base-corrected, and
refined branches using stripe-aware competitive attention.
In the simplified two-branch form,
\begin{equation}
  F_{\text{fuse}}
  =
  F_{\text{raw}} + w \odot (F_{\text{corr}} - F_{\text{raw}}).
  \label{eq:feature-fusion}
\end{equation}
The full framework generalizes this to three branches
$b \in \{1,2,3\}$ (raw, base-corrected, refined), where per-stripe fusion
weights are learned:
\begin{equation}
  F_{\text{fuse}}^{(s)}
  =
  \sum_{b=1}^{3} \alpha_b^{(s)} F_b^{(s)},
  \quad
  \alpha^{(s)}
  =
  \operatorname{softmax}\!\bigl(\operatorname{MLP}([\bar{F}_1^{(s)};\ldots;\bar{F}_3^{(s)}])/\tau\bigr),
\end{equation}
where $\bar{F}_b^{(s)}$ is the average-pooled descriptor of branch $b$ at
stripe $s$.
Equation~\eqref{eq:feature-fusion} is the special case $B{=}2$ with
$w{=}\alpha_{\text{corr}}^{(s)}$; in either form, the raw branch anchors the
representation so that correction acts as a bounded residual.
The gate also receives compact front-end statistics such as illumination
strength, color risk, rollback strength, and correction magnitude, so trust in
the corrected branch remains explicitly bounded by the state of the
illumination module.
This feature-space gate is the second trust mechanism: even if the image branch
is corrected aggressively, the representation need not trust it
unconditionally.

\subsection{Geometry-Guided Training}
\label{sec:geometry}

The framework is trained with three groups of objectives.
The first is the standard retrieval loss
$\mathcal{L}_{\text{reid}} = \mathcal{L}_{\text{cls}} +
\mathcal{L}_{\text{metric}}$,
instantiated with cross-entropy, triplet~\cite{hermans2017triplet},
ArcFace~\cite{deng2019arcface}, and center-loss terms~\cite{wen2016discriminative};
related metric-learning formulations such as FaceNet and Circle Loss provide
closely related geometric viewpoints~\cite{schroff2015facenet,sun2020circle}.
The second is a weak photometric feasibility term
$\mathcal{L}_{\text{photo}}$, which keeps the correction inside the bounded
set described above, together with an illumination-consistency term on
synthetic photometric variants,
$\mathcal{L}_{\text{iicl}} = V^{-1}\sum_{v}(1 - \cos(z, z^{(v)}))$.

The third group provides the key retrieval-oriented supervision beyond
standard descriptor-level ReID objectives~\cite{hermans2017triplet,deng2019arcface,dou2023isr,zhu2024seas}.
The training code exposes warmup, optional illumination-only, and joint
phases; the precise ATRW schedule is given in Section~\ref{sec:setup}.
After warmup, we freeze a snapshot of the raw-reference encoder as a teacher
and use it to supervise the geometry of corrected features.
A teacher-manifold tube loss constrains corrected features to remain inside a
same-identity neighborhood:
\begin{equation}
  \mathcal{L}_{\text{tube}}
  =
  [\|d^{\parallel}\| - r]_{+}^{2}
  +
  \lambda_{\perp}\|d^{\perp}\|_{2}^{2},
\end{equation}
where $d^{\parallel}$ and $d^{\perp}$ are projections onto the teacher
subspace and its complement.
We further use a teacher-manifold separation term
$\mathcal{L}_{\text{sep}}$ to maintain class boundaries and a differentiable
cross-light SoftAP objective $\mathcal{L}_{\text{softap}}$ to preserve
listwise retrieval ordering under illumination change.
To reduce mini-batch bias, the teacher gallery can be augmented with a queue of
frozen teacher descriptors from recent batches.
In the final ATRW mainline, other active auxiliaries are kept weak and
implementation-oriented: masked identity-image preservation and nuisance
decorrelation.
Together these terms specify not only whether correction helps, but which local
retrieval geometry the corrected feature should preserve.

The full objective is
\begin{equation}
  \mathcal{L}_{\text{total}}
  =
  \lambda_{\text{reid}}\mathcal{L}_{\text{reid}}
  +
  \lambda_{\text{photo}}\mathcal{L}_{\text{photo}}
  +
  \lambda_{\text{iicl}}\mathcal{L}_{\text{iicl}}
  +
  \lambda_{\text{geo}}\mathcal{L}_{\text{geo}}
  +
  \lambda_{\text{aux}}\mathcal{L}_{\text{aux}},
  \label{eq:total-loss}
\end{equation}
where
$\mathcal{L}_{\text{geo}} =
\mathcal{L}_{\text{tube}} + \mathcal{L}_{\text{sep}} + \mathcal{L}_{\text{softap}}$
and $\mathcal{L}_{\text{aux}}$ bundles the weak nuisance-decorrelation and
masked identity-image preservation terms.
During joint training, $\lambda_{\text{photo}}$ is annealed from a moderate
initial value to a small residual.
This schedule matches the theoretical view: photometric priors stabilize the
early search, then progressively yield to retrieval geometry.

\section{Experiments}

\subsection{Experimental Setup}
\label{sec:setup}

We evaluate primarily on ATRW~\cite{li2020atrw}, an open-set wildlife
re-identification benchmark, and report single-camera mAP, cross-camera mAP,
and their mean (mmAP) as the main selection metric.
We choose ATRW as the main testbed because it exposes the target problem most
clearly: severe cross-camera illumination variation coupled with
texture-dependent identity cues.
The experiments test three points: whether generic photometric correction alone
helps retrieval, whether trust-controlled adaptation recovers gains over that
generic correction, and whether the same ordering remains visible across
backbones, datasets, and human-oriented enhancement baselines.

\paragraph{Baseline.}
Unless otherwise stated, the main ablation uses
OSNet-AIN~\cite{zhou2021osnetain}.
We define a matched ReID baseline (hereafter: \emph{baseline}) as the same backbone, stripe head, input
resolution ($256{\times}384$), PK sampler ($P{=}8$, $K{=}4$), ReID loss
suite~\cite{luo2019bag,hermans2017triplet,deng2019arcface}, and total training
budget as the corresponding RIIC-ReID run, but without the illumination front-end
or geometric supervision.
For clarity, the matched main ablation contains four within-budget variants:
\emph{baseline}; \emph{naive illumination}, which keeps only a safe correction
branch; \emph{illumination only}, which adds trust-controlled correction but no
teacher geometry; and \emph{full RIIC-ReID}, which enables all components.
We additionally construct a generic illumination control that keeps a learned
safe-correction branch but removes model-aware residual refinement, identity
protection, stripe-wise rollback, raw/corrected branch fusion, and
teacher-guided geometry. This isolates whether standalone illumination
correction already helps retrieval, or only becomes useful after it is coupled
to internal encoder evidence and retrieval geometry.

\paragraph{Training schedule.}
The code retains a three-phase training interface, but the final ATRW mainline
collapses it into a practical two-stage run: a 15-epoch ReID warmup, a
zero-epoch standalone illumination phase, and 66 joint-training epochs.
Functionally, this is a warmup-plus-joint schedule with auxiliary weights
ramped during the early joint stage; the teacher encoder for
$\mathcal{L}_{\text{geo}}$ is frozen at the end of warmup.

\paragraph{Additional evidence.}
Beyond ATRW, we use standardized fixed query-gallery splits of GZGC Zebra and
StripeSpotter as auxiliary cross-species robustness checks, in line with the
broader move toward more explicit animal-ReID evaluation resources and
benchmarking protocols~\cite{cermak2024wildlifedatasets,adam2024seaturtleid,hou2025openanimals}.
These are not presented as official open-set benchmarks; final reported
numbers come from fixed query/gallery splits, while checkpoint selection uses
held-out train query/gallery splits.
Their role is to test whether the same model-centric ordering survives beyond
the main ATRW stress-test, not to claim a fully standardized universal suite
for cross-species re-identification.
We compare the same three variants as above: baseline, generic
illumination, and full model.
We also test whether human-oriented enhancement alone can explain the gain by
preprocessing ATRW and GZGC Zebra with a Zero-DCE++ implementation and
RetinexNet, then retraining the same OSNet-AIN baseline on the enhanced data.

\subsection{Main Results}

Table~\ref{tab:main-results} presents the primary comparison.
The pattern directly follows the logic of Propositions~1--3.
The naive illumination control underperforms the baseline, with
mmAP dropping from $56.69$ to $55.62$ and cross-camera mAP dropping from
$40.87$ to $39.60$.
This rules out the easy explanation that any learnable correction branch is
sufficient: without trust control, correction can perturb identity-relevant
structure on the hardest cross-camera metric.

Re-introducing task-aligned coupling reverses that behavior.
The task-aligned illumination front-end raises mmAP to $57.86$, which is
$+1.17$ over the baseline and $+2.24$ over the naive control.
This jump is the main evidence for trust-controlled adaptation: the problem is
not illumination correction per se, but uncontrolled correction.
Adding full teacher-guided geometry yields a smaller but still consistent
further improvement to $58.11$ mmAP, a
$+1.42$ gain over the baseline and a $+0.25$ gain over the
illumination-only front-end.
The improvement is again concentrated on cross-camera mAP
($41.80$ for RIIC-ReID vs.\ $40.87$ for the baseline and $39.60$ for the
naive control), which is the metric most exposed to illumination shift.

\begin{table}[t]
  \caption{ATRW main ablation under a matched protocol.
  Naive illumination removes model-centric trust and geometry coupling.}
  \label{tab:main-results}
  \centering
  \small
  \setlength{\tabcolsep}{4pt}
  \begin{tabularx}{\columnwidth}{@{}Y C{0.15\columnwidth} C{0.15\columnwidth} C{0.12\columnwidth} C{0.18\columnwidth}@{}}
    \toprule
    \thead[l]{Variant} & \thead{mAP\\$_{\text{single}}$} & \thead{mAP\\$_{\text{cross}}$} & mmAP & \thead{Rel.\\Baseline} \\
    \midrule
    Baseline & 72.52 & 40.87 & 56.69 & 0.00 \\
    Naive illumination & 71.65 & 39.60 & 55.62 & -1.07 \\
    Illumination only & 74.30 & 41.43 & 57.86 & +1.17 \\
    Full RIIC-ReID & \textbf{74.43} & \textbf{41.80} & \textbf{58.11} & \textbf{+1.42} \\
    \bottomrule
  \end{tabularx}
\end{table}

\subsection{Cross-Backbone Robustness}

The main table uses OSNet-AIN because it is the primary paper backbone, but
our claim is not that RIIC-ReID is tied to one encoder family.
To test this directly, we repeat the same baseline-versus-full-model comparison on
OSNet-X1.0 and ResNet-50 while keeping the backbone-specific setup fixed within
each row.
Table~\ref{tab:backbone-robustness} shows consistent positive gains:
$+1.42$ mmAP on OSNet-AIN, $+1.44$ on OSNet-X1.0, and $+1.01$ on ResNet-50,
with corresponding cross-camera mAP gains of $+0.93$, $+1.40$, and $+1.05$.
The benefit therefore does not come from one backbone's inductive bias alone.

\begin{table}[t]
  \caption{ATRW cross-backbone robustness under matched within-backbone
  protocols.}
  \label{tab:backbone-robustness}
  \centering
  \small
  \setlength{\tabcolsep}{3pt}
  \resizebox{\columnwidth}{!}{%
  \begin{tabular}{lcccccc}
    \toprule
    Backbone & Base mmAP & Full mmAP & $\Delta$mmAP & Base Cross & Full Cross & $\Delta$Cross \\
    \midrule
    OSNet-AIN X1.0 & 56.69 & 58.11 & +1.42 & 40.87 & 41.80 & +0.93 \\
    OSNet X1.0 & 59.06 & 60.50 & +1.44 & 43.35 & 44.75 & +1.40 \\
    ResNet-50 & 56.61 & 57.62 & +1.01 & 42.92 & 43.97 & +1.05 \\
    \bottomrule
  \end{tabular}%
  }
\end{table}

\subsection{Cross-Species Auxiliary Evidence}

We next ask whether the same ordering survives beyond ATRW on fixed
query-gallery evaluations with different species and capture conditions.
Table~\ref{tab:cross-species} reports Rank-1 and mAP on GZGC Zebra and
StripeSpotter.
On GZGC Zebra, generic illumination already gives a moderate gain
($72.71$ vs.\ $70.32$), but the full model improves further to $74.22$.
On StripeSpotter, generic illumination is slightly worse than the
baseline ($91.29$ vs.\ $91.55$), while the full model still reaches
$93.89$.
This asymmetry is informative: generic correction is not uniformly harmful, but
its effect is unstable, whereas the full task-aligned model remains best on
both datasets.

\begin{table}[t]
  \caption{Cross-species fixed query-gallery evidence on standardized splits.
  Model selection uses held-out train splits.}
  \label{tab:cross-species}
  \centering
  \small
  \setlength{\tabcolsep}{3pt}
  \begin{tabularx}{\columnwidth}{@{}L{0.19\columnwidth} Y C{0.12\columnwidth} C{0.12\columnwidth} C{0.16\columnwidth}@{}}
    \toprule
    \thead[l]{Dataset} & \thead[l]{Variant} & Rank-1 & mAP & \thead{Rel.\\Baseline} \\
    \midrule
    GZGC Zebra & Baseline & 71.70 & 70.32 & 0.00 \\
    GZGC Zebra & Generic illumination & 74.48 & 72.71 & +2.39 \\
    GZGC Zebra & Full model & \textbf{75.92} & \textbf{74.22} & \textbf{+3.90} \\
    \addlinespace[1pt]
    StripeSpotter & Baseline & 93.00 & 91.55 & 0.00 \\
    StripeSpotter & Generic illumination & 92.00 & 91.29 & -0.26 \\
    StripeSpotter & Full model & \textbf{94.00} & \textbf{93.89} & \textbf{+2.34} \\
    \bottomrule
  \end{tabularx}
\end{table}

\subsection{Generic Perceptual Enhancement Baselines}

We finally test whether the gain can be explained simply by better-looking
inputs produced outside the retrieval model.
Table~\ref{tab:perceptual-baselines} shows that on ATRW, Zero-DCE++ reaches
$57.75$ mmAP while RetinexNet stays at $56.67$; both remain below the full
model at $58.11$.
On GZGC Zebra, Zero-DCE++ and RetinexNet reach $67.31$ and $68.51$ mAP,
again well below the task-aligned full model at $74.22$.
Generic enhancement is therefore not useless, but it is weaker and less
reliable than retrieval-aligned correction.

\begin{table}[t]
  \caption{Generic human-perceptual enhancement baselines with a task-aligned
  full-model reference.
  Dashes denote metrics undefined under the respective protocol:
  ATRW's open-set evaluation does not define Rank-1 or overall mAP;
  GZGC~Zebra uses a single-partition gallery without the
  single-camera/cross-camera distinction.}
  \label{tab:perceptual-baselines}
  \centering
  \footnotesize
  \setlength{\tabcolsep}{2pt}
  \begin{tabular}{@{}L{0.17\columnwidth} L{0.16\columnwidth} C{0.075\columnwidth} C{0.075\columnwidth} C{0.11\columnwidth} C{0.11\columnwidth} C{0.095\columnwidth}@{}}
    \toprule
    \thead[l]{Dataset} & \thead[l]{Method} & Rank-1 & mAP & \thead{mAP\\$_{\text{single}}$} & \thead{mAP\\$_{\text{cross}}$} & mmAP \\
    \midrule
    ATRW & Zero-DCE++ & -- & -- & 73.14 & 42.36 & 57.75 \\
    ATRW & RetinexNet & -- & -- & 71.70 & 41.64 & 56.67 \\
    ATRW & Full RIIC-ReID & -- & -- & \textbf{74.43} & \textbf{41.80} & \textbf{58.11} \\
    \addlinespace[1pt]
    GZGC Zebra & Zero-DCE++ & 68.01 & 67.31 & -- & -- & -- \\
    GZGC Zebra & RetinexNet & 68.64 & 68.51 & -- & -- & -- \\
    GZGC Zebra & Full RIIC-ReID & \textbf{75.92} & \textbf{74.22} & -- & -- & -- \\
    \bottomrule
  \end{tabular}
\end{table}

\subsection{Interpretation}

Taken together, the results support a reliability claim rather than a universal
dominance claim. Under a wildlife-focused evaluation suite, the generic-control
row shows that learnable photometric correction is not automatically
retrieval-beneficial; on ATRW it hurts the hardest metric. Once trust control
and geometry are restored, bounded correction becomes useful and remains the
most stable variant across backbones, cross-species splits, and perceptual
baselines.

This is consistent with the logic of Section~3: the generic-control failure is the
empirical counterpart of Proposition~1, the weaker, less stable behavior of
generic enhancement is consistent with Proposition~2, and the gain from
rollback and branch fusion matches Proposition~3. The stronger claim is therefore
not that every non-task-aware corrector fails, but that generic correction is
unstable while retrieval-aligned correction is more reliable.

\section{Limitations}

Our formulation is intentionally conservative.
Section~3 gives a local feasible-manifold analysis under smoothness and local
curvature assumptions, not a proof of global optimality for deep retrieval
training.
Although the formulation is general, our strongest empirical validation is
currently concentrated on wildlife benchmarks, where illumination shift is a
salient nuisance; broader human-centered validation remains future work.
RIIC-ReID also remains most appropriate when illumination is a major nuisance and
identity depends on local texture structure; if lighting is not the bottleneck,
or if a much larger backbone can already absorb it, the utility of a structured
front-end may decrease.
In the converse case---controlled lighting with minimal illumination
variation---the rollback gates and trust mechanisms are designed to converge
toward the raw branch, but we have not empirically verified that the front-end
introduces zero residual perturbation in such benign settings.
Finally, RIIC-ReID is architecturally heavier than the baseline, and we do not
yet provide a dedicated runtime, parameter, or memory breakdown, so its benefit
must be judged against deployment cost.

\section{Conclusion}

We presented RIIC-ReID, a task-aligned illumination correction framework for
cross-light re-identification. The central claim is motivated by retrieval more
broadly: the
front-end should be optimized for retrieval geometry rather than perceptual
preference alone. We instantiate and primarily validate that view in wildlife
re-identification, where cross-camera illumination shift is especially severe.
Weak photometric priors define a feasible set, trust mechanisms regulate how
far correction is allowed to move, and teacher-guided geometry determines
which direction is worthwhile inside that set. Proposition~1 motivates task
alignment over pure photometric optimization, Proposition~2 motivates weak
photometric regularization, and Proposition~3 motivates rollback and
trust-bounded fusion.

Empirically, the contrast between RIIC-ReID and the generic illumination control
shows that illumination correction is not inherently beneficial for retrieval,
but becomes useful and reliable once it is constrained by feasibility,
moderated by trust, and selected by downstream retrieval structure.
Across ATRW, cross-backbone tests, standardized fixed query-gallery evaluations
on GZGC Zebra and StripeSpotter, and perceptual baselines such as Zero-DCE++
and RetinexNet, the ordering is consistent: the most reliable gains come from
learning illumination handling with the retrieval model rather than designing it
for human perception alone.

\bibliographystyle{ACM-Reference-Format}
\bibliography{acmmm_mm2026_refs}

\end{document}
